\section{Problem Setup}
\label{sec:problem}

We specify the cognitive state object and the formal problem addressed by the architecture.

\subsection{State object}
At time or session index $t$, the cognitive substrate is represented as a tuple
\begin{equation}
\label{eq:state-object}
  \ThoughtGraph_t = (V_t, E_t, \LayerSet, \ell, \tau, \mathrm{emb}, o, p),
\end{equation}
where $(V_t,E_t)$ is a directed typed graph, $\ell$ assigns layers, $\tau$ assigns edge types, $\mathrm{emb}$ maps vertices to embeddings when enabled, and $o,p$ record openness and persistence scope.
Equation~\eqref{eq:state-object} is \emph{provisional but close to exact}: additional bookkeeping (timestamps, agent provenance, audit hashes) may be appended without changing the core structure.

Let $\GlobalState_t$ denote the persistent global graph at $t$.
Let $\SessionState_t(q)$ denote session-local state for query $q$, constructed as a subgraph of $\GlobalState_t$ together with session-only vertices and edges for open thoughts.

\subsection{Retrieval and slicing}
Retrieval is modeled abstractly by an operator
\[
  \SessionState_t(q) = \SliceOp(q, \GlobalState_t),
\]
which may combine vector similarity on embeddings with relational neighborhood expansion.
The exact slice objective is an open design choice; the architecture only requires that slices be computable and auditable enough to support provenance (complexity or approximation guarantees for slicing are outside this paper's scope).

\subsection{Formal problem (architecture-level)}
The central design problem is to specify operators and constraints so that:
\begin{enumerate}[leftmargin=*]
  \item Session agents can explore freely using open thoughts without contaminating global state prematurely.
  \item Higher-layer closures respect intended support structure relative to lower layers.
  \item Layer $0$ remains externally grounded and is not silently overwritten by higher-layer operations.
\end{enumerate}

\begin{figure}[t]
  \centering
  % Figure 2 — Layered graph + support (PROJECT_SPEC.md 11A)
% Layout pass: left column for layer titles (wrapped); graph in open right region; semantics unchanged.
\resizebox{\linewidth}{!}{%
\begin{tikzpicture}[
  font=\small,
  nd/.style={circle, draw, minimum size=6mm, inner sep=0pt},
  layer/.style={draw, rounded corners, fill=gray!5, minimum width=9.0cm, minimum height=1.42cm},
  intra/.style={thick},
  sup/.style={-{Latex[length=2.2mm]}, thick, shorten >=1.2pt, shorten <=1.2pt},
  con/.style={-{Latex[length=2.2mm]}, thick, dashed, red!68!black, shorten >=1.8pt, shorten <=1.8pt},
]

\def\layersep{1.72}
\foreach \L/\lab in {0/Observations,1/Middle synthesis,2/Abstract beliefs} {
  \node[layer] (L\L) at (0,\L*\layersep) {};
  \node[anchor=west, text width=2.75cm, align=left, inner sep=0pt, xshift=7pt] at (L\L.west) {\textbf{Layer \L:} \lab};
}

% Graph shifted right so vertices sit in the clear band (not on layer titles)
\node[nd] (v0a) at (-1.15,0) {};
\node[nd] (v0b) at (0.95,0) {};
\node[nd] (v0c) at (3.15,0) {};
\draw[intra] (v0a) -- (v0b) -- (v0c);

\node[nd] (v1a) at (-0.05,\layersep) {};
\node[nd] (v1b) at (2.65,\layersep) {};
\draw[intra] (v1a) -- (v1b);

\node[nd] (v2a) at (1.15,2*\layersep) {};

\draw[sup] (v2a) -- node[right, font=\footnotesize, xshift=2pt, pos=0.42] {support} (v1a);
\draw[sup] (v2a) -- (v1b);
\draw[sup] (v1a) -- (v0b);
\draw[sup] (v1b) -- (v0c);

% Arc bulges upward, clear of the layer-1 title column
\draw[con] (v1a) to[bend left=28, looseness=1.12] (v1b);

\node[anchor=north west, font=\footnotesize, yshift=-0.34cm, xshift=0pt] at (L0.south west) {%
  \emph{Legend:} solid arrows, downward support/dependency along $\Tadm$; dashed red, contradiction (example).%
};

\end{tikzpicture}%
}

  \caption{Layered thought graph with downward support/dependency structure. Higher layers should admit dependency witnesses toward lower layers; contradiction and implication edges (if present) should remain visually and logically distinct. \emph{Simplification:} only one explicit support path is highlighted.}
  \label{fig:layers}
\end{figure}
